\section{The predicates and the modular bridge}\label{sec:games}
\subsection{Flagged magic square and its nine-copy gluing}
Let $\mathcal R=\{+++,+--,-+-,--+\}$ and let $\mathcal C_j$ be the triples of product $\eta_j$, where $\eta=(+,+,-)$. Put $r_0=c_0=+++$ and $r_1=c_1=+--$. Alice's questions are $R_1,R_2,R_3,T_0,T_1$; Bob's are $C_1,C_2,C_3,U_0,U_1$. A row answer is in $\mathcal R\cup\{f\}$, and a column answer in $\mathcal C_j\cup\{f\}$. On $T_x$ the answers are $\mathcal C_1\setminus\{c_x\}$ and two signs; on $U_y$ they are $\mathcal R\setminus\{r_y\}$ and two signs.

At $(R_i,C_j)$ accept two flags, reject a single flag, and otherwise require equal intersecting signs. At $(T_x,C_1)$ a sign matches the class $\{f,c_x\}$; a triple must match Bob's triple. At $(R_1,U_y)$ use the reversed rule with $r_y$. At $(T_x,U_y)$ two signs must satisfy $ab=(-1)^{xy}$; if either answer is a triple accept. All other cells accept all valid answers. Sample the 25 pairs uniformly. This is $\FMS$, in a real Pauli convention unitarily equivalent to \cite{FMS}.

For $\GIII$ take nine disjoint copies indexed by $j\in\Z/9\Z$ and one extra Bob question $\mathrm{Lab}$, with answers $a\in\Z/9\Z$. Besides the 225 internal pairs, use
\[
 (R_1^j,\mathrm{Lab}):\quad (\text{Alice answers }f)\iff(a=j),
\]
\[
 (R_2^j,\mathrm{Lab}):\quad (\text{Alice answers }+++)\iff(a\in\{j+1,j+2\}).
\]
These are the 243 equally weighted pairs. Other cross-copy pairs have weight zero. All indices below are modulo nine. The finite predicate is also enumerated in the ancillary data.

The real square used to define the ideal projectors is
\[
 O=\begin{pmatrix}X\otimes I&I\otimes X&X\otimes X\\
 I\otimes Z&Z\otimes I&Z\otimes Z\\
 X\otimes Z&Z\otimes X&Y\otimes Y\end{pmatrix}.
\]
Joint row and column eigenspaces give the four magic outcomes; append a one-dimensional flag. The calibrated rays are
\[
 c_0=(1,0,1,0)/\sq,\quad c_1=(0,1,0,1)/\sq,
\]
\[
 r_0=(1,1,1,1)/2,\quad r_1=(1,-1,1,-1)/2.
\]
Their overlap is $\langle c_x,r_y\rangle=(-1)^{xy}/\sq$. The sign projectors are the rays $(f\pm c_x)/\sq$ and $(f\pm r_y)/\sq$; all remaining outcomes are the calibrated residual magic rays. They win on
\[
 \Psi_2=(|ff\rangle+\sq\sum_{i=0}^3|ii\rangle)/3.
\]

\subsection{The conference-basis game}
For $\Gthree$, both players have base questions $R_1,R_2,R_3,C_1,C_2,C_3,H$. Each has four magic rays and a flag. The six old bases are the joint eigenspaces above. The additional orthonormal basis consists of the columns
\[
 r_0=\frac{(1,1,-1,0)^T}{\sqrt3},\quad
 r_1=\frac{(1,-1,0,1)^T}{\sqrt3},\quad
 r_2=\frac{(0,1,1,1)^T}{\sqrt3},\quad
 r_3=\frac{(1,0,1,-1)^T}{\sqrt3}.
\]
Their numerator Gram matrix is $3I$. Up to signed row and column permutations, these numerators form the classical skew conference matrix of order four, the $q=3$ case of the Paley construction; see Goethals--Seidel \cite[Section 2, Theorem 2.1]{Conference}. In the coordinate basis the $R_2$ labels are $d_0=+++$, $d_1=-+-$, $d_2=+--$, $d_3=--+$. Alice adds $T_0,T_1$, with signs and residual labels $k_j$ for $j\ne x$. Bob adds $U_0,U_1$, with signs and labels $h_j$ for $j\ne y$. Sample all 81 pairs uniformly and use the following rules.
\begin{enumerate}
\item On base/base cells accept two flags and reject one flag. Two magic outcomes are accepted exactly when $\Tr(P_{q,a}P_{r,b})>0$ for the specified rational ray projectors.
\item On $(T_x,R_2)$ a sign matches $\{f,d_x\}$ and $k_j$ matches $d_j$.
\item On $(H,U_y)$ a Bob sign matches $\{f,y\}$ and $h_j$ matches $j$.
\item On $(T_x,U_y)$ two signs obey $ab=(-1)^{xy}$; residual answers are accepted.
\item All remaining cells accept every valid pair.
\end{enumerate}
Thus the predicate uses rational zero tests only. There are five answers at every question and 732 forbidden events. Use $c_0=e_0,c_1=e_1$ and the displayed $r_0,r_1$ for the sign rays. The ideal state is
\[
 \Psi_3=(|ff\rangle+\sqrt3\sum_{i=0}^3|ii\rangle)/\sqrt{13}.
\]
Base orthogonality and calibration give zero forbidden amplitudes. For signs the amplitude is $(1+ab(-1)^{xy})/(2\sqrt{13})$, so the rejected parity amplitudes vanish.

\subsection{Extraction in arbitrary normal strategies}
\begin{lemma}[Checked classes]\label{lem:classes}
A perfectly matched coarse partition is a centralizer PVM, with its unique standard commutant mirror. Fine effects with sharp marginals are the corresponding joint spectral projections.
\end{lemma}
\begin{proof}
Use Lemma~\ref{lem:support}. For different classes $c,d$, the commuting positive product $E_cF_d$ has expectation zero and annihilates $\xi=\xi_\phi$. Summing gives $E_c\xi=E_cF_c\xi=F_c\xi$. Hence $E_cE_d\xi=E_cF_d\xi=0$. Separation and completeness give orthogonality and idempotence. A self-adjoint mirror gives $\phi(E_cx)=\phi(xE_c)$, so $E_c\in M_\phi$ and $F_c=JE_cJ$. A positive fine effect lies under each of its projection marginals; completeness saturates their products.
\end{proof}
Apply this to flag/magic classes and cell signs. On the magic projection $m=1-f$ the square relations hold internally. Write $a=O_{11},b=O_{12},c=O_{21},d=O_{22}$. They give
\[
 [a,b]=[a,c]=[b,d]=[c,d]=0,\qquad cb=-bc,\qquad ad=-da.
\]
Indeed the two expressions for $O_{33}$ are $acbd=-abcd$; multiplication by $a$ and $d$ gives the first anticommutation, and $[ac,bd]=0$ gives the second. Two commuting Pauli pairs therefore generate a unital centralizer $M_4$, with arbitrary multiplicity. No trace or finite-dimensional rigidity assumption has been used.

For $\Gthree$ the diagonal base cells check every outcome. An actual conference outcome $E_j$ is orthogonal to every old ray orthogonal to $r_j$: the forbidden event is $\phi(PE_jP)=0$ and faithfulness gives $E_jP=0$. The following old rays span the four required orthogonal complements:
\[
\begin{array}{c|ccc}
j&\multicolumn{3}{c}{\text{spanning rays}}\\\hline
0&e_3&e_0-e_1&e_0+e_2\\
1&e_2&e_0+e_1&e_0-e_3\\
2&e_0&e_1-e_2&e_1-e_3\\
3&e_1&e_0-e_2&e_0+e_3
\end{array}
\]
Consequently $E_j\le r_jr_j^T\otimes1$. Both four-tuples sum to $m$, so each positive difference is zero. This fixes the conference basis completely, including multiplicities.

\begin{lemma}[Modular bridge]\label{lem:bridge}
Every normal perfect strategy for $\FMS$, respectively $\Gthree$, on faithful support contains four partial isometries $\beta_i$ from $f$ onto the four magic corners, with modular eigenvalue $r=2$, respectively $r=3$. In particular
\[
 \phi(m)=4r\phi(f),\qquad \phi(f)=\frac1{1+4r}.
\]
\end{lemma}
\begin{proof}
Choose a minimal $q$ of the centralizer $M_4$ and centralizer maps $u_x,v_y$ from $q$ to the calibrated rays. Their products are $u_x^*v_y=(-1)^{xy}q/\sqrt r$. On the standard corner Hilbert space $eJeJ L^2(M)$, where $e=f+q$, pull back the sign differences through $f+u_x$ and $J(f+v_y)J$. They are commuting self-adjoint contractions $X_x,Y_y$. The joint-support vector is
\[
 \theta_{xy}=(f+u_x^*)J(f+v_y^*)J\xi
             =(f+(-1)^{xy}q/\sqrt r)\xi.
\]
The rejected sign products give
\[
 X_x\theta_{xy}=(-1)^{xy}Y_y\theta_{xy},\qquad
 X_xY_y\theta_{xy}=(-1)^{xy}\theta_{xy}.
\]
Each coefficient $f\pm q/\sqrt r$ is invertible in the centralizer of the corner. Its vector separates both corner algebras; cancellation first makes all four contractions reflections. Put $H=f+q/\sqrt r$ and $Z=f-q$. The three positive cells identify one $X$ and one $Y$ with $XH\xi=YH\xi$. In the negative cell, $XZH\xi=-YZH\xi=-ZXH\xi$, since $Y$ commutes with $Z$. Separation gives $XZ+ZX=0$.

It follows that $X=\beta+\beta^*$, with $\beta=qXf$, $\beta^*\beta=f$ and $\beta\beta^*=q$. The self-adjoint mirror makes $X$ centralize $\omega(x)=\phi(kx)$, $k=f+q/r$. The bounded invertible centralizer perturbation formula of Pedersen--Takesaki \cite[Lemma 4.5]{PT} is
\[
 \sigma_t^\omega(x)=k^{it}\sigma_t^\phi(x)k^{-it}.
\]
Its $q$--$f$ block gives $\sigma_t^\phi(\beta)=r^{it}\beta$. Multiplication by the magic matrix units gives all four maps. KMS gives $\phi(q)=r\phi(f)$. If one sector vanished, the off-diagonal reflection would kill the other, contradicting $f+m=1$. This proves nonvanishing and the mass formula.
\end{proof}
This is the arbitrary-normal extension of the flagged-overlap mechanism in \cite[Appendix C]{FMS}. It is not obtained by substituting type-II multiplicities into a finite-dimensional self-test.

\subsection{Semifinite exclusion}
For a perfect $\GIII$ strategy, the first links make the flags $F_j$ a partition of the identity, with Bob mirrors $JF_jJ$. The second links identify the designated magic corner with $C_j=F_{j+1}+F_{j+2}$. Lemma~\ref{lem:bridge} supplies
\begin{equation}\label{eq:ck}
 w_j^*w_j=F_j,\quad w_jw_j^*=C_j,\quad
 \sigma_t^\phi(w_j)=2^{it}w_j,\quad
 \sum_jF_j=1,\quad\sum_jC_j=2.
\end{equation}
The mass equations $2\phi(F_j)=\phi(F_{j+1})+\phi(F_{j+2})$ imply $\phi(F_j)=1/9$ by a maximum principle around the cycle. Every copy has full support, since its flag and magic masses total $9\phi(F_j)=1$.

Suppose $M$ is semifinite and $\phi=\Tr(h\,\cdot)$ on faithful support. Then $hw_j=2w_jh$ as affiliated operators. With $v_j=w_j^*$, bounded functional calculus gives $g(h)v_j=v_jg(h/2)$. For a bounded nonnegative $g$ compactly supported in $(0,\infty)$, trace integration and the two covers in \eqref{eq:ck} yield
\[
 \phi(g(h))=\sum_j\Tr(hg(h)v_jv_j^*)
 =\frac12\sum_j\Tr(C_jhg(h/2))=\phi(g(h/2)).
\]
For rigor use $\Tr(T^{1/2}vv^*T^{1/2})=\Tr(v^*Tv)$ for positive integrable $T$, rather than cycling arbitrary unbounded products. The probability law of $h$ under $\phi$ would have equal mass on all dyadic windows $(2^k,2^{k+1}]$. Every such mass is zero, a contradiction.

For a semifinite Bob ambient algebra, standardize Bob first. His standard commutant is semifinite and contains the compressed Alice effects, so the same proof applies. Conditioning on a positive-mass semifinite central summand preserves zero losing events; hence all state weight of either ambient algebra lies on its type-III central part. An unused semifinite summand need not be absent. More generally, if Alice's algebra is contained in a semifinite separator $N\subseteq B'$ and the state is normal on $N$, perfect play is impossible.
